%========================================================
% D3 — Correlation Tensor of the Phase-Evolved Bell State
%========================================================

\section{Correlation Tensor of the Phase-Evolved Bell State}

This derivation computes explicitly the correlation tensor associated with the reduced phase model.

\medskip

The phase-evolved Bell state is

\begin{equation}
|\psi(\theta)\rangle
=
\frac{1}{\sqrt{2}}
\left(
|01\rangle
+
e^{i\theta}|10\rangle
\right).
\label{eq:der_phase_bell_state}
\end{equation}

The corresponding density operator is

\begin{equation}
\rho(\theta)
=
|\psi(\theta)\rangle\langle\psi(\theta)|.
\label{eq:der_phase_density_operator}
\end{equation}

Expanding explicitly gives

\begin{align}
\rho(\theta)
&=
\frac{1}{2}
\left(
|01\rangle
+
e^{i\theta}|10\rangle
\right)
\left(
\langle 01|
+
e^{-i\theta}\langle 10|
\right)
\nonumber
\\
&=
\frac{1}{2}
\Big(
|01\rangle\langle 01|
+
e^{-i\theta}|01\rangle\langle 10|
+
e^{i\theta}|10\rangle\langle 01|
+
|10\rangle\langle 10|
\Big).
\label{eq:der_phase_density_expansion}
\end{align}

\medskip

The correlation tensor is defined by

\begin{equation}
T_{ij}(\theta)
=
\mathrm{Tr}
\left[
\rho(\theta)
\left(
\sigma_i \otimes \sigma_j
\right)
\right],
\qquad
i,j \in \{x,y,z\}.
\label{eq:der_correlation_tensor_definition}
\end{equation}

Equivalently, since the state is pure,

\begin{equation}
T_{ij}(\theta)
=
\langle \psi(\theta) |
\sigma_i \otimes \sigma_j
|
\psi(\theta)
\rangle.
\label{eq:der_correlation_tensor_pure_state}
\end{equation}

The Pauli matrices act on the computational basis as

\begin{align}
\sigma_x |0\rangle &= |1\rangle,
&
\sigma_x |1\rangle &= |0\rangle,
\nonumber
\\
\sigma_y |0\rangle &= i|1\rangle,
&
\sigma_y |1\rangle &= -i|0\rangle,
\nonumber
\\
\sigma_z |0\rangle &= |0\rangle,
&
\sigma_z |1\rangle &= -|1\rangle.
\label{eq:der_pauli_basis_action}
\end{align}

\medskip

We now compute each nonzero tensor component.

\subsection*{Diagonal components}

For \( T_{xx} \), one has

\[
(\sigma_x \otimes \sigma_x)|01\rangle = |10\rangle,
\qquad
(\sigma_x \otimes \sigma_x)|10\rangle = |01\rangle.
\]

Therefore,

\begin{align}
T_{xx}(\theta)
&=
\frac{1}{2}
\left(
\langle 01|
+
e^{-i\theta}\langle 10|
\right)
\left(
|10\rangle
+
e^{i\theta}|01\rangle
\right)
\nonumber
\\
&=
\frac{1}{2}
\left(
e^{i\theta}
+
e^{-i\theta}
\right)
\nonumber
\\
&=
\cos\theta.
\label{eq:der_Txx_phase_model}
\end{align}

For \( T_{yy} \), one has

\[
(\sigma_y \otimes \sigma_y)|01\rangle = |10\rangle,
\qquad
(\sigma_y \otimes \sigma_y)|10\rangle = |01\rangle.
\]

Hence the same calculation gives

\begin{equation}
T_{yy}(\theta)
=
\cos\theta.
\label{eq:der_Tyy_phase_model}
\end{equation}

For \( T_{zz} \), one has

\[
(\sigma_z \otimes \sigma_z)|01\rangle = -|01\rangle,
\qquad
(\sigma_z \otimes \sigma_z)|10\rangle = -|10\rangle.
\]

Thus,

\begin{align}
T_{zz}(\theta)
&=
\langle \psi(\theta) |
\sigma_z \otimes \sigma_z
|
\psi(\theta)
\rangle
\nonumber
\\
&=
-1.
\label{eq:der_Tzz_phase_model}
\end{align}

\medskip

\subsection*{Off-diagonal transverse components}

For \( T_{xy} \), one has

\[
(\sigma_x \otimes \sigma_y)|01\rangle
=
-i|10\rangle,
\qquad
(\sigma_x \otimes \sigma_y)|10\rangle
=
i|01\rangle.
\]

Therefore,

\begin{align}
T_{xy}(\theta)
&=
\frac{1}{2}
\left(
\langle 01|
+
e^{-i\theta}\langle 10|
\right)
\left(
-i|10\rangle
+
i e^{i\theta}|01\rangle
\right)
\nonumber
\\
&=
\frac{1}{2}
\left(
i e^{i\theta}
-
i e^{-i\theta}
\right)
\nonumber
\\
&=
-\sin\theta.
\label{eq:der_Txy_phase_model}
\end{align}

For \( T_{yx} \), one has

\[
(\sigma_y \otimes \sigma_x)|01\rangle
=
i|10\rangle,
\qquad
(\sigma_y \otimes \sigma_x)|10\rangle
=
-i|01\rangle.
\]

Therefore,

\begin{align}
T_{yx}(\theta)
&=
\frac{1}{2}
\left(
\langle 01|
+
e^{-i\theta}\langle 10|
\right)
\left(
i|10\rangle
-
i e^{i\theta}|01\rangle
\right)
\nonumber
\\
&=
\frac{1}{2}
\left(
-i e^{i\theta}
+
i e^{-i\theta}
\right)
\nonumber
\\
&=
\sin\theta.
\label{eq:der_Tyx_phase_model}
\end{align}

\medskip

\subsection*{Vanishing components}

All remaining components involving one transverse Pauli operator and one \( z \)-Pauli operator vanish:

\begin{equation}
T_{xz}
=
T_{zx}
=
T_{yz}
=
T_{zy}
=
0.
\label{eq:der_vanishing_phase_tensor_components}
\end{equation}

This follows because these operators map the subspace spanned by \( |01\rangle \) and \( |10\rangle \) to the orthogonal subspace spanned by \( |00\rangle \) and \( |11\rangle \), or produce contributions with zero overlap against \( |\psi(\theta)\rangle \).

\medskip

Collecting the computed components gives

\begin{equation}
T(\theta)
=
\begin{pmatrix}
T_{xx} & T_{xy} & T_{xz}\\
T_{yx} & T_{yy} & T_{yz}\\
T_{zx} & T_{zy} & T_{zz}
\end{pmatrix}
=
\begin{pmatrix}
\cos\theta & -\sin\theta & 0\\
\sin\theta & \cos\theta & 0\\
0 & 0 & -1
\end{pmatrix}.
\label{eq:der_correlation_tensor_phase_model}
\end{equation}

\medskip

The result shows that the reduced phase \( \theta \) acts as a rotation of the transverse correlation components in the \( x\text{-}y \) plane, while the longitudinal anticorrelation \( T_{zz}=-1 \) remains unchanged.

This is consistent with the physical interpretation of the reduced phase model: the local phase shift modifies coherence between \( |01\rangle \) and \( |10\rangle \), but does not alter their populations.